% Einheit 3 – Normalform und p-q-Formel
\clearpage
\setcounter{aufgabe}{24}
\hypertarget{e3}{}\einheitenkopf{Einheit 3 von 4 · Normalform und p-q-Formel}
\zweigzeile{Hier lernst du, eine quadratische Gleichung in die Normalform zu bringen und mit der p-q-Formel zu lösen · neu in diesem Jahr (an der Oberschule je nach Buch Klasse 9 oder 10) · P10 oft · baut auf: Wurzelziehen und Lösbarkeit (Einheit 1), Satz vom Nullprodukt (Einheit 2), Terme ordnen, binomische Formeln}

\begin{aufgabe}{Ich erkenne, welche Form eine quadratische Gleichung hat und welcher Lösungsweg passt. Kreuze die Form an und schreibe den Lösungsweg dazu: Wurzelziehen, Nullprodukt, Rückwärtsrechnen oder Formel. Löse nichts.}
\sachtabelle{lccccc}{Gleichung & \shortstack{ohne\\$x$-Glied} & \shortstack{Produkt\\$= 0$} & \shortstack{Klammer$^2$\\$=$ Zahl} & \shortstack{Normal-\\form} & Lösungsweg}{a) \ $x^2 = 70$ & $\square$ & $\square$ & $\square$ & $\square$ & \leerzelle\leerzelle \\ b) \ $(x - 2)\cdot(x + 9) = 0$ & $\square$ & $\square$ & $\square$ & $\square$ & \leerzelle\leerzelle \\ c) \ $(x + 1)^2 = 7$ & $\square$ & $\square$ & $\square$ & $\square$ & \leerzelle\leerzelle \\ d) \ $x^2 - 3x + 1 = 0$ & $\square$ & $\square$ & $\square$ & $\square$ & \leerzelle\leerzelle \\ e) \ $3x^2 - 27 = 0$ & $\square$ & $\square$ & $\square$ & $\square$ & \leerzelle\leerzelle}
\end{aufgabe}

\begin{aufgabe}{Ich kann $p$ und $q$ mit Vorzeichen ablesen. Eine Gleichung in Normalform hat die Gestalt $x^2 + px + q = 0$. Schreibe $p$ und $q$ mit Vorzeichen auf. Rechne nichts.}
\begin{geruest}
\gz{a}{$x^2 + 5x + 4 = 0$}{\feld{p}\feld{q}}
\gz{b}{$x^2 - 3x + 2 = 0$}{\feld{p}\feld{q}}
\gz{c}{$x^2 + x - 12 = 0$}{\feld{p}\feld{q}}
\gz{d}{$x^2 - 11x = 0$}{\feld{p}\feld{q}}
\gz{e}{$x^2 - 2 = 0$}{\feld{p}\feld{q}}
\end{geruest}
\end{aufgabe}

\begin{aufgabe}{Ich erkenne, ob die Vorzahl von $x^2$ eins ist. Kreuze an, ob du vor der Formel teilen musst, und gib an, durch welche Zahl. Rechne nichts.}
\begin{teile}
\teil $2x^2 + 6x - 8 = 0$ \hfill \kreuz{teilen}\kreuz{nicht teilen}\quad teilen durch: \leerfeld
\teil $x^2 + 4x - 1 = 0$ \hfill \kreuz{teilen}\kreuz{nicht teilen}\quad teilen durch: \leerfeld
\teil $-x^2 + 3x + 4 = 0$ \hfill \kreuz{teilen}\kreuz{nicht teilen}\quad teilen durch: \leerfeld
\teil $\frac{1}{2}x^2 + x - 3 = 0$ \hfill \kreuz{teilen}\kreuz{nicht teilen}\quad teilen durch: \leerfeld
\teil $5x^2 - x = 0$ \hfill \kreuz{teilen}\kreuz{nicht teilen}\quad teilen durch: \leerfeld
\end{teile}
\end{aufgabe}

\begin{aufgabe}{Ich kann eine Gleichung in Normalform mit der p-q-Formel lösen. Lies $p$ und $q$ ab und setze sie in die Formel $x_{1,2} = -\frac{p}{2} \pm \sqrt{\left(\frac{p}{2}\right)^2 - q}$ ein.}
\beispiel{x^2 + 6x + 8 &= 0 && p = 6,\ q = 8 \\ x_{1,2} &= -3 \pm \sqrt{9 - 8} = -3 \pm 1 \\ x_1 &= -2, \quad x_2 = -4}
\begin{gleichungsraster}[2]
\gl{x^2 + 6x + 5 = 0} & \gl{x^2 + 10x + 21 = 0} \\ \noalign{\vspace{6pt}}
\gl{x^2 + 12x + 35 = 0} & \gl{x^2 + 14x + 40 = 0} \\ \noalign{\vspace{6pt}}
\gl{x^2 - 4x - 5 = 0} & \\ \noalign{\vspace{6pt}}
\end{gleichungsraster}
\end{aufgabe}

\begin{aufgabe}{Ich kann eine Gleichung in Normalform mit der p-q-Formel lösen – weiter. Bringe die Gleichung, wenn nötig, zuerst in die Normalform: alles auf eine Seite, rechts null, Vorzahl von $x^2$ eins. Löse dann mit der Formel.}
\begin{gleichungsraster}[3]
\gl{x^2 = 2x + 24} & \gl{-2x^2 + 12x - 10 = 0} \\ \noalign{\vspace{6pt}}
\gl[2]{x^2 - 5x + 6 = 0} & \gl[2]{x^2 - 10x + 25 = 0} \\ \noalign{\vspace{6pt}}
\gl[2]{x^2 + 2x + 7 = 0} & \\ \noalign{\vspace{6pt}}
\end{gleichungsraster}
\medskip
\begin{teile}
\teil Der Wert unter der Wurzel, $D = \left(\frac{p}{2}\right)^2 - q$, heißt Diskriminante. Berechne $D$ für $x^2 + 7x + 11 = 0$ und gib an, wie viele Lösungen die Gleichung hat, ohne sie zu lösen.
\end{teile}
\end{aufgabe}

\begin{aufgabe}{Ich kann eine Gleichung in Normalform mit der p-q-Formel lösen – weiter. Die Wurzel geht nicht auf: Gib die Lösungen mit Wurzel und als Näherungswerte auf zwei Nachkommastellen an.}
\begin{gleichungsraster}[2]
\gl{x^2 - 4x + 1 = 0} & \\ \noalign{\vspace{6pt}}
\end{gleichungsraster}
\medskip
\begin{teile}
\teil Prüfe durch Einsetzen, ob $x = -7$ eine Lösung von $x^2 + 5x - 14 = 0$ ist. \hfill \kreuz{(wA)}\kreuz{(fA)}
\end{teile}
\medskip
Löse zuerst die Klammer auf und ordne. Löse dann mit der Formel.
\begin{gleichungsraster}[3]
\gl{(x + 3)^2 = 2x + 14} & \\ \noalign{\vspace{6pt}}
\end{gleichungsraster}
\end{aufgabe}

\begin{aufgabe}{Ich kann unterscheiden, welcher Lösungsweg zu einer Gleichung passt. Löse jede Gleichung auf dem kürzesten Weg: Wurzelziehen, Nullprodukt, Rückwärtsrechnen oder Formel. Schreibe den Weg dazu.}
\begin{gleichungsraster}[2]
\gl{4x^2 = 100} & \gl{(x - 6)\cdot(x + 1) = 0} \\ \noalign{\vspace{6pt}}
\gl{x^2 - 9x = 0} & \gl{(x - 4)^2 = 1} \\ \noalign{\vspace{6pt}}
\gl{x^2 + x - 6 = 0} & \\ \noalign{\vspace{6pt}}
\end{gleichungsraster}
\medskip
\begin{teile}
\teil Berechne die Nullstellen der Funktion $f$ mit $f(x) = 2x^2 + 4x - 4$. Runde auf zwei Nachkommastellen. (P10 2020)
\end{teile}
\end{aufgabe}

\begin{aufgabe}{Ich kann die Schnittpunkte einer Parabel und einer Geraden berechnen. Setze die beiden Terme gleich, löse die Gleichung und berechne die $y$-Werte der Schnittpunkte.}
\begin{teile}
\teil Parabel $y = x^2 - 3$ \ und \ Gerade $y = 2x + 5$
\teil Parabel $y = (x - 3)^2 - 2$ \ und \ Gerade $y = -x + 3$ \ (P10 2022)
\end{teile}
\end{aufgabe}

\begin{aufgabe}{Ich finde den Fehler in einer Rechnung mit der p-q-Formel. Sara löst $x^2 - 14x + 45 = 0$ so:}
\rechnung{x^2 - 14x + 45 &= 0 && p = -14,\ q = 45 \\ x_{1,2} &= -7 \pm \sqrt{49 - 45} = -7 \pm 2 \\ x_1 &= -5, \quad x_2 = -9}
\begin{teile}
\teil Finde den Fehler, benenne ihn und korrigiere die Rechnung.
\end{teile}
\medskip
Löse selbst.
\begin{gleichungsraster}[2]
\gl{x^2 - 8x + 7 = 0} & \\ \noalign{\vspace{6pt}}
\end{gleichungsraster}
\end{aufgabe}

\begin{aufgabe}{Ich kann begründen, wann die p-q-Formel gilt und wie viele Lösungen es gibt.}
\begin{teile}
\teil Begründe, warum man bei $3x^2 - 6x - 9 = 0$ zuerst durch $3$ teilen muss, bevor man die p-q-Formel anwendet.
\teil Begründe: Ist $q$ negativ, dann hat $x^2 + px + q = 0$ immer zwei Lösungen.
\end{teile}
\end{aufgabe}

\begin{aufgabe}{Ich kann berechnen, wo ein geworfener Ball auf dem Boden aufkommt. Ein Ball fliegt auf der Bahn $y = -0{,}1x^2 + 0{,}6x + 1{,}6$. Dabei ist $x$ die waagerechte Entfernung vom Abwurfpunkt und $y$ die Höhe über dem Boden, beide in Metern.}
\begin{teile}
\teil In welcher Entfernung kommt der Ball auf dem Boden auf? Welche Lösung passt nicht zur Frage?
\teil In welcher Entfernung ist der Ball genau $2\,\text{m}$ hoch? Runde auf zwei Nachkommastellen.
\end{teile}
\end{aufgabe}
